We next construct the core measure independently. Normalize a common primary model \(Y\), with finite normalization \(N\to Y\). Use only its top normal components; lower components and their boundary points remain in the domain but have zero degree-\(n\) contribution. Let \(Z\) run through their exceptional prime divisors of dimension \(p\). These are finitely many integral proper \(k\)-schemes, with DVR generic rings. Put \(b_Z=\operatorname{ord}_Z(\mathfrak a)>0\), \(\nu_Z=\operatorname{ord}_Z/b_Z\in X\). The valuation is pulled back from the fraction field of its minimal component, with that component as support. For negative-effective real curve-nef core divisors \(D_1,\ldots,D_p\), define
\[
\mu_Y^{\rm geom}(D_1,\ldots,D_p)
=\sum_Z b_Z(D_1|Z\cdots D_p|Z\cdot Z)\delta_{\nu_Z}. \tag{A}
\]
Pullbacks to \(N\) are understood. For \(p=0\), the intersection is \(\deg_k Z=[\kappa(Z):k]\).

Positivity precedes ideal comparison. G331's qualified real nef-plus-ample and finite-span rationalization give negative rational ample divisors approaching each \(D_i\) in coefficient space, by subtracting small \(E_I\)-multiples and rounding negative-effective coefficients upward. Finite normalization pulls ample bundles to ample bundles; restriction to \(Z\) stays ample. Clearing rational denominators and complete 33.45.9 give positive mixed numerical intersections. Their finite coefficient polynomials converge, so the weights in (A) are nonnegative. This argument also covers arbitrary real nef entries and needs no definition of ideal positivity.

For an exceptional real Cartier \(F\) on this model, G331's evaluation at \(\nu_Z\) is \(\operatorname{ord}_Z(F)/b_Z\). Normal divisor cycles and n3's projection yield the geometric identity
\[
\int_X H_F\,d\mu_Y^{\rm geom}=Q_Y(D_1,\ldots,D_p,F). \tag{G}
\]
Here \(H_E=v(I)\) for an effective \(E_I\); negative divisors evaluate negatively.

The atomic measure is invariant under refinements independently of algebraic moments. If \(Y'\to Y\) is a primary refinement, its normalization \(N'\) maps to \(N\): on each dominant normal component, integral functions defining \(N\) belong to the integrally closed local rings of \(N'\). This lift is proper (use the closed graph over \(Y\)) and birational on top components. For a prime \(Z'\) mapping dominantly to a dimension-\(p\) prime \(Z\) of \(N\), its DVR dominates \(\mathcal O_{N,\eta_Z}\) in the SAME fraction field. They are equal: every field element is a unit times a power of the lower uniformizer, and its sign of order is preserved. Their orders, residue fields, \(b\)'s and valuations therefore coincide. Apply cycle projection (P) to the pulled-back effective primary \(E_I\), whose order at \(Z\) is positive. It forces exactly one such prime over each \(Z\), with degree one. Numerical projection gives its mixed weight equal to the old weight. Every other \(Z'\) has image dimension below \(p\); its \(p\) pulled-back line bundles have intersection zero by 33.45.7. Thus (A) is invariant as a measure on \(X\). The same proof handles common refinements and auxiliary normal lifts, with every top component considered separately. In particular (G) can now be tested against Cartier functions determined on ANY common refinement.

We prove the bridge to n2's algebraic measures using actual image ideals, not only valuation identities. Fix \(Y=\operatorname{Bl}_I R\). Qualified G331 gives \(T=H^0(Y,\mathcal O_Y)\) finite birational, \(T/R\) supported at \(\mathfrak m\). Choose \(h\) with \(\mathfrak m^h(T/R)=0\) and fix ONCE
\(\mathfrak c_0=\mathfrak m^{h+1}\), so \(\mathfrak c_0T\subset\mathfrak m\) as elements of \(R\). For a negative rational ample exceptional \(L\) and sufficiently large divisible \(q\), \(\mathcal O_Y(qL)\) is globally generated. Set
\[
F_q=H^0(Y,\mathcal O_Y(qL))\subset T,\qquad J_q=\mathfrak c_0F_q\subset R.
\]
This is an ideal, proper because \(\mathfrak c_0T\subset\mathfrak m\), and unit off \(\mathfrak m\) by localization of global sections and triviality of the divisor there. Hence it is proper primary.

Fix \(Y'=\operatorname{Bl}_{I\mathfrak c_0}R\), dominating \(Y\) and \(\operatorname{Bl}_{\mathfrak c_0}R\); write \(D_0\) for the negative divisor of the latter. Global generation means equality of the IMAGE subsheaf of \(F_q\otimes\mathcal O_Y\) in \(\mathcal O_Y\) with the invertible \(\mathcal O_Y(qL)\). Pullback preserves its surjection. The pullback inclusion into \(\mathcal O_{Y'}\) remains injective, since its Cartier equation stays regular on every reduced dominating component. Multiplying by the invertible image of \(\mathfrak c_0\) proves the exact ideal identity
\[
J_q\mathcal O_{Y'}=\mathcal O_{Y'}(qf^*L+D_0). \tag{S}
\]
Consequently the ACTUAL blowup universal property gives \(Y'\to\operatorname{Bl}_{J_q}R\). This is stronger than equality of valuation functions.

For simultaneous \(L_1,\ldots,L_n\) and exponents \(q_i\), n3's (I) and projection give
\[
\frac{e(J_{1,q_1},\ldots,J_{n,q_n})}{q_1\cdots q_n}
=-Q_{Y'}(f^*L_1+D_0/q_1,\ldots,f^*L_n+D_0/q_n). \tag{S2}
\]
On every support, (S) also gives
\[
-v(J_q)/q=H_L-v(\mathfrak c_0)/q. \tag{S3}
\]
The error is uniform because \(\mathfrak c_0\) is PRIMARY. The right side of (S2) converges by the finite-dimensional Cartier polynomial on ONE fixed model.

For real negative curve-nef \(D_i\), take negative rational ample \(L_{i,j}\to D_i\) in their fixed finite coefficient spans, as in the positivity argument, and choose generated divisible \(q_{i,j}\to\infty\). The correction and refinement are unchanged. Both (S2)'s Cartier polynomials and (S3)'s function profiles converge. N2's uniform moment continuity, including its mass bound for varying continuous tests, gives
\[
\int_X H_{D_n}\,d\mu^{\rm alg}(H_{D_1},\ldots,H_{D_p})
=Q_Y(D_1,\ldots,D_p,D_n). \tag{B}
\]
Thus no uniform-function-to-intersection convergence was assumed.

For arbitrary integral exceptional \(F\), choose \(N\) so \(F-NE_I\) and \(-NE_I\) are both negative ample, using G331's Cartier order unit and ample twisting. Their difference is \(F\); linearity extends (B) to this test and then to every real exceptional Cartier test. A primary ideal test is determined on a common primary refinement. The independent atomic invariance proved above, (G), and (B) identify all its moments on that refinement. G331's dense ideal-difference lattice therefore gives
\[
\mu_Y^{\rm geom}(D_1,\ldots,D_p)
=\mu^{\rm alg}(H_{D_1},\ldots,H_{D_p}). \tag{COMPARE}
\]
This proves independence from all core presentations giving the same functions, as well as exact geometric Cartier comparison and primary-model independence. The algebraic extension of n2 is consequently the bounded product extending this independently constructed geometric core.
